\documentclass[UTF8,a4paper,11pt]{ctexart}
\usepackage[top=2.1cm,bottom=2.1cm,left=2.3cm,right=2.3cm]{geometry}
\usepackage{amsmath,amssymb,booktabs,tabularx,array,xcolor,setspace,enumitem,fancyhdr,xurl,hyperref}
\usepackage[ruled,vlined,linesnumbered]{algorithm2e}
\definecolor{ink}{HTML}{263449}
\hypersetup{colorlinks=true,linkcolor=ink,urlcolor=ink,
 pdftitle={核心算法与复杂度附录：当前实现与适用边界},
 pdfauthor={西安交通大学；陈若凡及团队}}
\setstretch{1.15}
\setlength{\emergencystretch}{2em}
\setlength{\parskip}{3pt}
\setlength{\headheight}{15pt}
\pagestyle{fancy}
\fancyhf{}
\fancyhead[L]{\small XH-202631\quad 核心算法与复杂度附录}
\fancyhead[R]{\small 西安交通大学}
\fancyfoot[C]{\thepage}
\renewcommand{\headrulewidth}{0.3pt}
\setlist[itemize]{leftmargin=1.6em,itemsep=2pt,topsep=3pt}
\SetAlgorithmName{算法}{算法}{算法目录}
\SetKwInput{KwIn}{输入}
\SetKwInput{KwOut}{输出}
\renewcommand{\thealgocf}{A\arabic{algocf}}
\newcommand{\codefile}[1]{{\footnotesize\path{#1}}}
\newenvironment{compactalgorithm}{\begingroup\small\setstretch{1.06}\begin{algorithm}[H]}{\end{algorithm}\endgroup}
\newcommand{\topic}[1]{\par\noindent\textbf{#1}\quad}
\begin{document}
\thispagestyle{plain}
\begin{center}
{\LARGE\bfseries\color{ink}核心算法与复杂度附录}\par
\vspace{4pt}
{\small 面向复杂任务求解的多智能体动态协作与自治交付系统}\par
{\small 西安交通大学\quad 代码核对日期：2026 年 9 月 14 日}
\end{center}

本附录给出五项核心机制的工程化伪代码、代码入口及复杂度边界。算法独立编号为 A1--A5，分别对应主报告第三至第六章；不新增实验结果。伪代码是多个实现函数的摘要，不是可直接执行的程序。符号与保守上界统一见第\ref{sec:complexity}节。

\section{需求前沿驱动的版本化任务图扩展}
\topic{作用}将尚未完成且依赖已满足的需求转化为下一阶段子图，保留已有节点的执行结果；阶段数目标本身不作为“任务成功”的依据。

\begin{compactalgorithm}
\caption{需求前沿选择、阶段规划与受控合并}
\KwIn{当前图 $G$，需求进度 $F$，全局任务契约 $T$，阶段预算 $b$}
\KwOut{扩展图 $G^{+}$，或保留原图并记录阻塞原因}
从 $G$ 刷新需求进度并持久化 $F$\;
\If{存在失败或阻塞的需求生产节点}{记录暂停原因；\Return 不扩展\;}
\If{不存在待规划需求}{按需追加一次全局治理节点；\Return 当前或治理后的图\;}
根据需求依赖、批量上限和提示预算选择就绪需求集合 $F_b$\;
\If{$F_b=\varnothing$}{记录需求依赖阻塞；\Return 不扩展\;}
调用阶段规划器生成候选；编译为阶段 TaskGraph $\Delta G$\;
为新增节点和内部引用添加阶段命名空间\;
检查物理产物写入归属、全局产物契约及代码节点限制\;
将新增根节点接到原图中非治理叶节点之后\;
构造 $G^{+}$：保留原节点状态，追加新节点，递增图版本\;
检查节点 ID、依赖存在性、DAG、能力与产物所有权\;
\eIf{阶段规划或合法性检查失败}{
 保留已完成节点；记录可恢复的阶段阻塞及失败原因\;
}{
 保存 $G^{+}$、需求生产者映射和阶段事件\;
}
\end{compactalgorithm}

\topic{实现入口}
\codefile{orchestration/core/workflow.py}：\texttt{expand\_requirement\_frontier}；
\codefile{orchestration/planning/staged_expansion.py}：\texttt{merge\_stage\_graph}、\texttt{append\_governance\_nodes}；
\codefile{orchestration/graph/task_graph.py}：\texttt{validate\_graph}。

\topic{不变量与边界}扩展不清空已完成节点，不允许新增节点越过全局物理产物写入契约。当前是阶段前沿扩展，不宣称能在任意执行指令处并发改图。规划异常可留下阻塞证据，但不会因此直接产生成功交付。

\topic{复杂度口径}合并后的图含 $n$ 个节点、$m$ 条依赖边。节点复制和依赖连接属于控制面开销，完整扩展还包含需求处理、契约检查和全图合法性检查，不能仅按新增节点数计费；模型生成时间另计。回归入口：\codefile{tests/test_requirement_staging.py}。

\clearpage
\section{稀疏依赖上下文与受保护记忆构造}
\topic{作用}依赖消息与运行期记忆走两条互补路径：前者只读取直接依赖，后者按相关性检索历史记录；目标、约束和结构化证据不交给有损摘要任意改写。

\begin{compactalgorithm}
\caption{直接依赖消息与 ContextCapsule 的组合构造}
\KwIn{图 $G$，当前节点 $v$，任务契约 $T$，运行记忆 $\mathcal M$，通信与记忆预算}
\KwOut{节点执行上下文，或上下文构造异常}
从图中取得规范节点与其直接依赖结果\;
\ForEach{直接依赖结果}{
 投影到消息字段白名单，保留声明的必需字段\;
 对普通文本限额压缩；对大结构化值或重复内容生成带哈希的内容引用\;
}
检查单消息及节点依赖上下文预算，必要时进一步外置内容\;
\If{受保护字段或引用元数据仍无法满足硬预算}{抛出预算异常，禁止交付该上下文\;}
读取本轮有效记忆，分离核心目标约束与普通候选\;
构建节点索引，逐条计算相关性、时效、重要性、图邻近度和证据质量\;
排序普通候选，取配置允许的前 $k$ 条\;
压缩所选普通摘要；将全局目标、关键约束、图位置和引用独立放入 capsule\;
\If{目标为空或关键约束数量校验不通过}{抛出受保护字段异常\;}
保存 capsule 与通信指标，组合为执行器输入\;
\end{compactalgorithm}

\topic{实现入口}
\codefile{orchestration/communication/context_builder.py}：\texttt{ContextBuilder}；
\codefile{orchestration/memory/runtime_memory.py}：\texttt{build\_capsule}、\texttt{\_score}；
\codefile{orchestration/graph/graph_scheduler.py}：\texttt{\_context\_for}。

\topic{接口边界}依赖消息预算与记忆摘要预算分别实施。完整模型提示还由
\codefile{orchestration/core/prompt_builder.py} 组装和检查，不能仅凭某一摘要满足预算，就宣称整个请求一定满足模型窗口。内容引用有哈希校验；当前 capsule 保存为结构化 JSON，不将所有字段都描述成统一签名对象。

\topic{复杂度}设有效记忆条数为 $R$、记忆总内容规模为 $S_R$，一次评分涉及的祖先遍历工作量为 $a_v$，其最坏上界为 $O(n+m)$。候选检索、评分与排序的控制开销可保守写为：
\[
T_{\mathrm{memory}}
=O\!\left(n+S_R+R\,a_v+R\log(R+1)\right).
\]
SQLite 查询、有损压缩和上下文序列化另计；节点索引、已加载记录和排序数组都会占用内存，不能把空间仅写成输出的 $k$ 条记录。仅当单条评分和记录长度有界时，候选处理才可简化为 $O(R\log R)$。

\topic{可核查产物}依赖内容引用、通信指标及 \texttt{memory/capsules/} 下的节点上下文。测试入口：\codefile{tests/test_communication.py}、\codefile{tests/test_runtime_memory.py}、\codefile{tests/test_prompt_budget_semantics.py}。

\clearpage
\section{受保护任务契约下的 Checkpoint 恢复}
\topic{作用}在已持久化的执行窗口边界恢复任务图和进度。完整性校验、目标一致性与完成状态共同约束恢复，而非重新解释用户需求。

\begin{compactalgorithm}
\caption{窗口级 checkpoint 的验证与继续执行}
\KwIn{运行目录，checkpoint 引用，预期目标和 TaskSpec}
\KwOut{恢复后的任务图及窗口状态，或拒绝恢复}
检查显式引用不是绝对路径且不含上级目录跳转；未指定时选最新 checkpoint\;
读取 JSON；重新计算去除 integrity 字段后的规范化内容哈希\;
\If{文件缺失、摘要不符或缺少完整任务图}{拒绝恢复并报告原因\;}
以 TaskGraph 模型重新验证图结构和契约字段\;
\If{保护目标与图目标不一致}{拒绝恢复\;}
\If{调用方提供预期目标或 TaskSpec}{
 对比目标、必需产物列表与受保护任务契约摘要\;
 \If{任一比较失败}{拒绝恢复，不自动复用不一致的旧状态\;}
}
恢复图与 epoch，登记已完成节点集合并记录恢复事件\;
按当前依赖状态选择后续 ready 窗口，已完成节点不重新进入该队列\;
窗口结束后保存任务图、节点状态、受保护上下文和新内容哈希\;
由正常执行与验收流程决定是否完成，而非仅检查 epoch 数量\;
\end{compactalgorithm}

\topic{实现入口}
\codefile{orchestration/core/workflow.py}。主要函数：
\codefile{load_horizon_checkpoint}、
\codefile{_execute_graph_in_horizons}、\codefile{_protected_task_digest}。

\topic{保护范围}主流程恢复时传入预期目标与任务契约。当前契约摘要覆盖任务 ID、任务类型、必需产物、成功条件和报告章节等字段，不应称为对整个 TaskSpec 任意字段的全量校验。该入口也不逐一重算所有外部 artifact 的哈希；产物检查由执行与验收路径分别承担。

\topic{恢复保证}已落盘的 completed 节点可以在继续执行时跳过；这不是任意崩溃时刻的 exactly-once 保证，也不恢复浏览器会话、外部模型内部状态或全部内存对象。SHA256 是内容完整性摘要，不是带身份认证的数字签名。

\topic{复杂度}令 checkpoint 序列化字节数为 $S_K$，图检查成本为 $T_{\mathrm{graph}}$。读取、解析和哈希的基础开销为 $O(S_K)$，恢复还需加上 $T_{\mathrm{graph}}$。当前规范化 JSON 使用键排序，其额外比较工作量另记为 $T_{\mathrm{canon}}$，不将其无条件忽略。未指定引用时还需枚举并排序 $K$ 个 checkpoint，成本至多 $O(K\log(K+1))$。内存至少覆盖解析后的快照规模。

\topic{业务长程的独立实现}城市千步实验使用
\codefile{task_plugins/urban_multimodal/long_horizon.py} 中的 \texttt{WorkUnitLedger.resume\_at}，校验 checkpoint、账本长度、目标摘要和末记录摘要；完整链验证另由业务验证器完成。该实现与通用任务图恢复不可混为一个入口。测试入口：\codefile{tests/test_long_horizon_continuity.py}、\codefile{tests/test_urban_long_horizon.py}。

\clearpage
\section{能力与资源约束下的动态执行器路由}
\topic{作用}在同能力执行器之间进行可解释选择，把运行历史反馈到排序中；失败后的重试、切换或重规划受恢复策略约束，而不是由模型自由决定。

\begin{compactalgorithm}
\caption{约束过滤、历史感知排序与有界恢复}
\KwIn{TaskNode $v$，执行器注册表，资源注册表，历史统计，已尝试集合与恢复预算}
\KwOut{执行器与 RoutingDecision，或无候选／恢复终止状态}
扫描注册执行器，保留 available、能力匹配且未被排除的候选\;
调用资源过滤函数检查位置、位置可用性、安全等级与配置的质量阈值等条件\;
按资源绑定和条件分支检查时延限制；移除不通过候选\;
\If{候选为空}{抛出无可用执行器异常，由控制面处理\;}
读取候选历史；依据配置采用静态排序或历史感知 effective score\;
按评分及稳定次序排列候选，选择首项\;
保存候选排名、评分分量、执行器／资源标识和排除集合\;
将执行器绑定到当前节点，构造上下文并执行\;
记录成功、质量、Token 和耗时等反馈\;
\If{节点执行失败}{
 依据错误类型和策略，选择同执行器重试、未尝试执行器切换、重规划或失败\;
 不增加 TaskSpec 允许的恢复预算；上下文构造失败不走普通执行器切换\;
}
将业务结果交给独立验证路径；执行返回成功不等于最终验收成功\;
\end{compactalgorithm}

\topic{实现入口}
\codefile{orchestration/execution/capability_registry.py}：\codefile{select_executor}、\codefile{_resource_allowed}；
\codefile{orchestration/execution/executor_scoring.py}；
\codefile{orchestration/graph/recovery_controller.py}：\texttt{decide\_after\_failure}。

\topic{当前评分}历史感知策略采用静态先验与运行历史的渐进混合：
\[
s(e)=(1-\lambda_e)s_{\mathrm{static}}(e)
     +\lambda_e s_{\mathrm{runtime}}(e)-p_{\mathrm{failure}}(e),
\qquad
\lambda_e=w_{\mathrm{history}}\frac{N_e}{N_e+c}.
\]
其中 $N_e$ 是历史运行次数，$c$ 为冷启动参数；分量来自质量、成本、时延、成功率和 Token 效率等。该式对应当前历史感知评分器；静态策略则保留质量、成本、时延和执行器标识的稳定排序。

\topic{复杂度}令 $M$ 为注册执行器数，$C$ 为资源条目数。当前资源过滤可能为每个执行器重复扫描注册资源，因此选择成本的保守上界为：
\[
T_{\mathrm{route}}=O\!\left(MC+M\log(M+1)\right),
\qquad S_{\mathrm{route}}=O(M+C).
\]
上述空间指候选与资源遍历临时结构，不含历史库。仅在 $C$ 有界或资源检查已索引化时，才可把主要选择成本简化为 $O(M\log M)$。

\topic{边界}端、边、云为逻辑资源后端；位置可用性检查不能替代物理资源逐项健康实测。当前决策记录包含合法候选等信息，不宣称已完整记录所有被过滤资源的逐项淘汰原因。测试入口：\codefile{tests/test_dynamic_routing.py}、\codefile{tests/test_resource_routing.py}。

\clearpage
\section{验证驱动的责任定位与局部重放}
\topic{作用}把独立验证发现的问题定位到责任生产节点，重开责任节点及其后继，并保留未受影响节点的结果。业务修复预算不作为普通执行器重试预算的隐式延长。

\begin{compactalgorithm}
\caption{证据归因、受影响节点重开与有界再验证}
\KwIn{当前图 $G$，业务验证结果，证据核验问题，最大修复轮数 $q$}
\KwOut{无需修复、已修复、failed 或 blocked}
从业务检查与证据核验结果构造带责任节点的问题集合 $I$\;
记录阻断问题；若验证通过且 $I=\varnothing$，返回无需修复\;
\If{存在已耗尽执行恢复的普通失败节点}{返回 failed，不额外授予业务重试\;}
\For{$j\leftarrow 1$ \KwTo $q$}{
 从 $I$ 提取可证明的责任节点集合 $O$\;
 \If{$O=\varnothing$}{记录无法路由的问题；返回 blocked\;}
 构造修复上下文；搜索 $O$ 及其传递后继 $V'$\;
 重开 $V'$ 中节点，清除旧执行结果与执行器绑定；保留其他节点状态\;
 登记本轮目标、影响范围和问题 ID，调用正常调度器重放\;
 \If{出现不可继续业务修复的执行失败}{记录本轮失败；返回 failed\;}
 重新运行领域验证，结合证据节点结果形成新问题集合 $I^{+}$\;
 依据验证证据关闭已消失问题，保存本轮修复结果\;
 \If{验证通过且 $I^{+}=\varnothing$}{返回已修复，继续后续需求验收与最终门禁\;}
 \If{连续两轮无进展或预算耗尽}{返回 blocked\;}
 $I\leftarrow I^{+}$\;
}
\end{compactalgorithm}

\topic{实现入口}
\codefile{orchestration/graph/repair_loop.py}：\texttt{BusinessRepairLoop.run}；
\codefile{orchestration/graph/task_graph.py}：\texttt{descendants\_of}、\texttt{reopen\_subgraph}；
\codefile{orchestration/core/workflow.py}：修复后的全局验收衔接。

\topic{范围与不变量}当前采用传递后继闭包作为保守影响范围，不宣称这是最小语义影响子图。范围之外的已完成结果不因本轮修复被清空；范围之内必须重新产生可验证结果。修复函数返回 repaired 后，仍需后续全局验收，不能直接替代 Final Gate。

\topic{复杂度}当前 \texttt{descendants\_of} 反复扫描全图，而不是从预先维护的后继邻接表直接遍历。设实际扫描轮数为 $h\leq n+1$，则影响定位和重开节点的控制开销为：
\[
T_{\mathrm{affected}}
=O\!\left(h(n+m)+n\,n'\right)+T_{\mathrm{topo}},
\]
其中 $n'=|V'|$，$n\,n'$ 来自逐个重开节点时的线性查找；$T_{\mathrm{topo}}$ 是全图拓扑排序成本。

若每轮还包含重放 $T_{\mathrm{replay},j}$、验证 $T_{\mathrm{validate},j}$ 和归因／记录 $T_{\mathrm{record},j}$，总时间必须按轮次求和。仅优化图搜索并不能消除模型和业务执行成本。测试入口：\codefile{tests/test_business_repair_loop.py}、\codefile{tests/test_business_validation.py}。

\clearpage
\section{复杂度分析的符号、假设与图控制面}\label{sec:complexity}
\topic{统一计量}以下为代码结构推导的保守上界，不是实测耗时或紧确界。字典与集合操作按通常的期望常数时间计；节点标识、单个配置字段和单条历史摘要按长度有界计。可变长正文、JSON、artifact 必须另按字节规模计入；本地存储 I/O、模型调用及外部工具不藏入常数项。

\begin{center}
\small
\begin{tabularx}{\textwidth}{@{}p{0.24\textwidth}X@{}}
\toprule
符号 & 含义\\
\midrule
$n,m$；$n',m'$ & 全图节点／边数；受影响子图节点／边数\\
$M,C$ & 注册执行器数；资源注册表条目数（两者不混用）\\
$R,k,S_R$ & 有效记忆条数；选取条数；记忆总内容规模\\
$a_v$ & 单条记忆评分中的图邻近度遍历工作量\\
$S_K,K$ & checkpoint 序列化规模；checkpoint 文件数\\
$S_A,S_{\mathrm{msg}}$ & 被哈希 artifact 字节数；消息实际内容规模\\
$h,h_b,q$ & 影响传播扫描轮数；失败阻塞传播轮数；修复轮数上限\\
$T_{\mathrm{topo}},T_{\mathrm{canon}}$ & 当前拓扑排序成本；JSON 规范化键排序成本\\
\bottomrule
\end{tabularx}
\end{center}

\subsection{图合法性检查：线性理论界的适用条件}
在邻接表、常数时间节点索引及 FIFO 队列下，标准 DAG 检查为 $O(n+m)$。当前 \texttt{validate\_graph} 对每个节点 ID 调用列表 \texttt{count}，仅重复检查就需 $O(n^2)$；\texttt{topological\_order} 还使用 \texttt{pop(0)}、候选队列重排序及子节点排序。因此，对当前实现采用：
\[
T_{\mathrm{topo}}
=O\!\left(n^2\log(n+1)+m\log(n+1)\right)
\]
作为保守上界。完整合法性检查还需计入输入输出契约、路径与 Schema 验证的内容成本。拓扑辅助结构为 $O(n+m)$；这一区分保留了标准算法理论界与当前确定性排序实现之间的差别。

\subsection{就绪扫描：不只扫描节点，还要解析依赖}
\texttt{get\_node} 通过线性列表查找；一次全图依赖检查可达到 $O(nm+n)$。\texttt{ready\_nodes} 先传播失败阻塞，再查询就绪状态并排序。在一次调用内，若阻塞传播经历 $h_b$ 轮，则保守成本为：
\[
T_{\mathrm{ready}}
=O\!\left((h_b+1)(nm+n)+n\log(n+1)\right).
\]
因此，完整调度应按实际调用次数累加该成本，并计入状态保存与执行开销，不能仅因“每轮扫描所有节点”就统一写成 $O(n^2+m)$。窗口化执行减少活动范围和部分持久化次数，但不自动消除这些全图查询。

\subsection{受影响子图：执行范围小不等于定位开销只与子图有关}
当前实现的影响传播与排序读取全图，A5 已给出相应成本。只有维护好后继索引、采用访问集合遍历，并使用常数时间节点定位时，局部搜索／重开的基础部分才可达到 $O(n'+m')$；若还要求稳定拓扑顺序，仍需另计排序成本。这些是可优化方向，不是本次提交宣称已经完成的性能保证。

\clearpage
\section{其他模块、端到端成本与复核边界}
\begin{center}
\small
\renewcommand{\arraystretch}{1.18}
\begin{tabularx}{\textwidth}{@{}p{0.18\textwidth}p{0.34\textwidth}X@{}}
\toprule
模块 & 当前计量口径 & 必要限定\\
\midrule
资源路由 & $O(MC+M\log(M+1))$ & 资源位置与可用性查询可能重复扫描；不把资源数与执行器数混为一个变量。\\
记忆构造 & 扫描、逐条评分、全候选排序，加查询与压缩 & 只写 $O(R)$ 仅覆盖扫描，遗漏排序、文本处理和图邻近度计算。\\
checkpoint 写入 & $O(S_K)+T_{\mathrm{canon}}$ & JSON 编码、键排序、哈希及文件 I/O 分别计量。快照包含的结果越大，$S_K$ 越大。\\
checkpoint 恢复 & 读取与规范化校验，加全图检查；必要时加文件排序 & 不把账本加载、完整图检查或业务重新验证当作常数时间。\\
artifact 哈希 & 时间 $O(S_A)$ & 分块流式哈希的额外缓冲可有界；一次性读入或 JSON 重编码实现则需额外内容空间。\\
稀疏依赖消息 & 按实际依赖数与 $S_{\mathrm{msg}}$ 计量 & 哈希、序列化、必需字段投影及预算调整都有内容成本；边数不等于字节数或 Token 数。\\
业务候选验证 & 按候选数与领域验证器计量 & 对 $P$ 个对象直接进行两两 AABB 检测，基础比较为 $O(P^2)$；不代表整个装箱求解器的总复杂度。\\
\bottomrule
\end{tabularx}
\end{center}

\topic{持久化成本}执行器统计按执行器／任务类型保存。若历史条目数为 $J$、序列化规模为 $S_J$，全量排序保存需额外考虑 $O(J\log(J+1)+S_J)$。图快照、事件日志、上下文和 artifact 的实际写入次数也应分别统计；此开销不包含在单次候选排序中。

\topic{端到端时间}对于实际发生的调度、恢复和验收序列，应采用：
\[
\begin{aligned}
T_{\mathrm{run}}
={}&\sum T_{\mathrm{control}}+\sum T_{\mathrm{model}}+\sum T_{\mathrm{tool}}\\
 &+\sum T_{\mathrm{validation}}+\sum T_{\mathrm{I/O}}.
\end{aligned}
\]
各项按实际执行路径计数，避免将模型时延解释为图算法性能，或将局部修复次数与全部运行成本混为一个指标。本附录不新增速度提升比例、不将千个 WorkUnit 等同于千次模型调用，也不宣称物理端边云集群已完成实测。

\topic{代码与证据复核}A1--A5 均给出源码路径与测试入口；代码在提交系统源码包中，历史结果在运行证据中。普通回归测试与历史真实模型运行是不同证据类型；测试文件存在本身不是“本次全部测试通过”的声明。复杂度推导以本次核对的具体数据结构为准，不将标准算法的理想界直接替代工程实现界。

\topic{改进方向}节点 ID 索引、后继邻接表、增量 ready 队列、资源可用性索引与候选 Top-$k$ 选择可分别降低查询、传播和排序成本。实施这些优化后应重新验证状态语义和验收结果，再更新复杂度，不以修改文档替代实现与测试。
\end{document}
